\documentclass[10pt,a4paper]{amsart}
\usepackage[margin=19mm]{geometry}
\usepackage{amsmath,amssymb,mathtools}
\usepackage[scr=rsfs,cal=euler]{mathalfa}
\usepackage{xfrac}
\usepackage[unicode,hidelinks,bookmarks=false]{hyperref}
\newcommand{\ie}{i.e.\ }
\newcommand{\cf}{cf.\ }
\newcommand{\resp}{respectively\ }
\newcommand{\Subsection}{\subsection}
\providecommand{\nobreakdash}{\nobreak\hbox{-}}
\setlength{\emergencystretch}{2em}
\multlinegap0pt
\usepackage{bm}
\usepackage{bbm}
\usepackage{dsfont}
\usepackage{xspace}
\usepackage{cancel}
\usepackage{mathtools}
\usepackage{amsthm}
\makeatletter
\@ifundefined{theorem}{\newtheorem{theorem}{Theorem}}{}
\@ifundefined{definition}{\newtheorem{definition}[theorem]{Definition}}{}
\@ifundefined{lemma}{\newtheorem{lemma}[theorem]{Lemma}}{}
\makeatother
\makeatletter
\newcommand{\C}{\ensuremath{\mathbb{C}}}
\expandafter\ifx\csname Corollary\endcsname\relax
\newenvironment{Corollary}{ \begin{corollary}\normalfont\rmfamily}{\end{corollary}}
\fi
\expandafter\ifx\csname Definition\endcsname\relax
\newenvironment{Definition}{ \begin{definition}\normalfont\slshape}{\end{definition}}
\fi
\expandafter\ifx\csname Example\endcsname\relax
\newenvironment{Example}{ \begin{example}\normalfont\rmfamily}{\bend\end{example}}
\fi
\expandafter\ifx\csname Lemma\endcsname\relax
\newenvironment{Lemma}{ \begin{lemma}\normalfont\slshape}{\end{lemma}}
\fi
\expandafter\ifx\csname Problem\endcsname\relax
\newenvironment{Problem}{ \begin{problem}\normalfont\slshape}{\end{problem}}
\fi
\expandafter\ifx\csname Proposition\endcsname\relax
\newenvironment{Proposition}{ \begin{proposition}\normalfont\rmfamily}{\end{proposition}}
\fi
\newcommand{\R}{\ensuremath{\mathbb{R}}}
\expandafter\ifx\csname Remark\endcsname\relax
\newenvironment{Remark}[1][noisnotdefined]{ \ifthenelse{\equal{#1}{noisnotdefined}}{\begin{remark}}{\begin{remark}[#1]}\normalfont\rmfamily}{\bend\end{remark}}
\fi
\expandafter\ifx\csname Theorem\endcsname\relax
\newenvironment{Theorem}[1][noisnotdefined]{ \ifthenelse{\equal{#1}{noisnotdefined}}{\begin{theorem}}{\begin{theorem}[#1]}\normalfont\slshape}{\end{theorem}}
\fi
\renewcommand{\d}{\, \mathrm{d}}
\newcommand{\dx}{\mathrm{d}}
\newcommand{\e}{\mathrm{e}}
\newcommand{\im}{\mathrm{i}}
\renewcommand{\vec}[1]{\ensuremath \mathbf{\boldsymbol{#1}}}
\newcommand{\w}{\ensuremath{\vec \omega}}
\newcommand{\x}{\ensuremath{\vec x}}
\makeatother
\title[ANOVA-boosting for Random Fourier Features]{ANOVA-boosting for Random Fourier Features}
\author{Daniel Potts, Laura Weidensager}
\date{Source: arXiv:2404.03050v2 (CC BY 4.0)}
\begin{document}
\maketitle

\noindent\emph{Source and licence.} ANOVA-boosting for Random Fourier Features. Version arXiv:2404.03050v2 (CC BY 4.0), distributed under the Creative Commons Attribution 4.0 International licence, as recorded in the arXiv licence metadata for this exact version and shown on the abstract page https://arxiv.org/abs/2404.03050v2. Full rights notice: licenses/arxiv-2404.03050-CC-BY-4.0.txt.\par\smallskip

\noindent\textit{Editorial scope.} Counted unit: the Lemma whose source label is \texttt{lem:ANOVA-terms\_f} in the article (source0.tex lines 458--467, source label \texttt{lem:ANOVA-terms\_f}), delivered together with the article's own complete original proof of it (source0.tex lines 468--503, one \texttt{proof} environment of 36 lines). No mathematics is written, repaired or re-derived: statement and proof are the article's own text line for line. Required context retained uncounted: def:anova-terms (lines 430--450); eq:mu\_prod (lines 432--434); eq:anova-terms (lines 436--439); eq:anova-decomp (lines 441--443). Every definition, display and label of those lines is retained unchanged. Omitted from this package and not redistributed: 1--429, 435--435, 440--440, 444--457, 504--2309. Nothing inside a retained excerpt is omitted or altered: every retained body is delivered exactly as the article's lines are written, so no omission receipt of any kind accompanies this package. Citations: the retained excerpts carry no citation command at all, so no bibliography entry of the article is transcribed and no reference is redistributed. Reference adaptation: none was needed and none was made; every reference inside the delivered unit resolves inside it, so no unresolved reference or citation is produced. Printed slips kept literal: no printed word, symbol, hypothesis or number of the counted statement or of its proof was altered. Layout and class: the article's own class is \texttt{article} with packages including arxiv, hyperref, amsfonts, amsmath, amssymb, amsthm, algorithm, algorithmic, graphicx, color, subcaption, mathtools, numprint, booktabs; the wrapper is the lane's pinned \texttt{amsart} editorial wrapper at 10pt on A4, which restates the theorem-like environments the retained text needs and adds the article's own macro definitions that the retained text needs (\texttt{\textbackslash C}, \texttt{\textbackslash R}, \texttt{\textbackslash d}, \texttt{\textbackslash dx}, \texttt{\textbackslash e}, \texttt{\textbackslash im}, \texttt{\textbackslash vec}, \texttt{\textbackslash w}, \texttt{\textbackslash x}), with extra wrapper packages (bm, bbm, dsfont, xspace, cancel, mathtools). The delivered numbering is the unit's own. Assets: no retained span contains a figure environment, a graphics-inclusion command, a PDF-page inclusion, a table or a TikZ picture; no image file is copied into this package, the package declares no asset dependency and no image file is redistributed with it. Licence: version arXiv:2404.03050v2 of this article is distributed under the Creative Commons Attribution 4.0 International licence, as recorded in the arXiv licence metadata for this exact version and shown on the abstract page https://arxiv.org/abs/2404.03050v2. A full rights notice is delivered at licenses/arxiv-2404.03050-CC-BY-4.0.txt. Characters: every retained excerpt is pure ASCII, so the delivered engine, XeLaTeX with its default Latin Modern text font, reports no missing character.
\begin{definition}\label{def:anova-terms}
Let $f$ be in $L_2(\R^d,\mu)$. For a tensor product measure 
\begin{equation}\label{eq:mu_prod}
\mu(\vec x) = \prod_{i=1}^d \mu_i(\vec x_i)
\end{equation}
and for a subset $\vec u\subseteq[d]$ we define the \textbf{ANOVA (Analysis of variance) terms} recursively by 
\begin{align}\label{eq:anova-terms}
f_{\varnothing} &= \int_{\R^d} f(\x) \mu(\x) \dx \x \notag\\
f_{\vec u}(\vec x_{\vec u}) &=\int_{\R^{d-|\vec u|}}f(\vec x) \mu_{\vec u^c}(\vec x_{\vec u^c})\d \vec x_{\vec u^c}-\sum_{\vec v\subset \vec u}f_{\vec v}(\vec x_{\vec v}).
\end{align}
The\textbf{ ANOVA decomposition} with respect to $\mu$ of a function ${f\colon \R^d\to \C}$ is then given by
\begin{equation}\label{eq:anova-decomp}
f(\vec x)=f_{\varnothing}+\sum_{i=1}^d f_{\{i\}}(x_i)+\sum_{i\neq j=1}^d f_{\{i,j\}}(x_i,x_j)+\cdots+f_{[d]}(\vec x)=\sum_{\vec u\subseteq[d]}f_{\vec u}(\vec x_{\vec u}).
\end{equation}
The terms~\eqref{eq:anova-terms} are the unique decomposition~\eqref{eq:anova-decomp}, 
such that 
\begin{align}
\langle f_{\vec u},f_{\vec v} \rangle_\mu \coloneqq \int_{\R^d} f_{\vec u}(\x_{\vec u})f_{\vec v}(\x_{\vec u})\mu(\x) \dx \x &= 0 \quad \text{ for } \vec v\neq \vec u\subseteq[d],\label{eq:orth}\\
\int_{\R} f_{\vec u}(\x_{\vec u}) \mu_j(x_j)\dx x_j &= 0 \quad \text{ for } j\in  \vec u. \label{eq:zero_mean}
\end{align}
\end{definition} 

\begin{lemma}\label{lem:ANOVA-terms_f}
Let the sampling distribution $\mu$ have a product structure~\eqref{eq:mu_prod}. Then the ANOVA decomposition \eqref{eq:anova-terms} for a function $f\in L_2(\R^d,\mu) \cap L_2(\R^d)$, is given by
\begin{equation}\label{eq:anova_hat}
f_{\vec u}(\vec x_{\vec u}) 
= \frac{1}{(2\pi)^d}\int_{\R^d} \hat f(\w)E(\x,\w,\mu, \vec u)   \d \w.
\end{equation}



\end{lemma}

\begin{proof}
The Fourier transform of the tensor product density $\mu$ can be decomposed as 
\begin{equation*}
\hat{\mu} (\w) = \prod_{i\in [d]} \hat{\mu}_i(\omega_i).
\end{equation*} 
We prove~\eqref{eq:anova_hat} inductively over $|\vec u|$. First, observe that
\begin{align*}
f_\varnothing &= \int_{\R^d} f(\vec x) \mu(\x) \d \vec x 
= \int_{\R^d}\frac{1}{(2\pi)^d}\int_{\R^d}\hat f(\w)\, \e^{\im \langle \w,\vec x\rangle} \d \w\,\mu(\x) \d \vec x\\
&=\frac{1}{(2\pi)^d}\int_{\R^d}\hat f(\w)\int_{\R^d}\, \e^{\im \langle \w,\vec x\rangle} \mu(\x) \d \x\d \w 
=\frac{1}{(2\pi)^d}\int_{\R^d}\hat f(\w) \hat{\mu}(-\w)\d \w \\
&= \frac{1}{(2\pi)^d}\int_{\R^d}\hat f(\w) E(\x,\w,\mu, \varnothing)\d \w,
\end{align*}
where $E(\x,\w,\mu, \varnothing) = \prod_{i\in [d]} \hat{\mu}_i(-\omega_i)$ does not depend on $\x$.
The induction step follows using Definition~\ref{def:anova-terms},
\begin{align*}
&f_{\vec u}(\vec x_{\vec u})
=\int_{\R^{d-|\vec u|}}\frac{1}{(2\pi)^d}\int_{\R^d}\hat f(\w)\, \e^{\im \langle \w,\vec x\rangle} \d \w \mu_{\vec u^c}(\vec x_{\vec u^c})\d \vec x_{\vec u^c}-\sum_{\vec v\subset \vec u}f_{\vec v}(\vec x_{\vec v})\\
&\quad=\frac{1}{(2\pi)^d}\int_{\R^d}\hat f(\w)\e^{\im \langle \w_{\vec u},\vec x_{\vec u}\rangle}\int_{\R^{d-|\vec u|}}\e^{\im \langle \w_{\vec u^c},\vec x_{\vec u^c}\rangle}  \mu_{\vec u^c}(\vec x_{\vec u^c})\d \vec x_{\vec u^c}\d \w
-\frac{1}{(2\pi)^d}\sum_{\vec v\subset \vec u}\int_{\R^d} \hat f(\w)E(\x,\w,\mu, \vec v)\d \w\\
&\quad=\frac{1}{(2\pi)^d}\int_{\R^d}\hat f(\w)\e^{\im \langle \w_{\vec u},\vec x_{\vec u}\rangle} \hat{\mu}_{\vec u^c}(-\w_{\vec u^c})\d \w-\frac{1}{(2\pi)^d}\sum_{\vec v\subset \vec u}\int_{\R^d} \hat f(\w) \prod_{i\in \vec v} \left(\e^{\im \omega_i x_i}- \hat{\mu}_i(-\omega_i)\right) \prod_{i \in \vec v^c} \hat{\mu}_i(-\omega_i) \d \w\\
&\quad=\frac{1}{(2\pi)^d}\int_{\R^d}\hat f(\w)\hat{\mu}_{\vec u^c}(-\w_{\vec u^c}) \left(\e^{\im \langle \w_{\vec u},\vec x_{\vec u}\rangle}-\sum_{\vec v\subset \vec u} \prod_{i\in \vec v} \left(\e^{\im \omega_i x_i}- \hat{\mu}_i(-\omega_i)\right) \prod_{i \in \vec u\backslash v} \hat{\mu}_i(-\omega_i)\right) \d \w\\
&\quad=\frac{1}{(2\pi)^d}\int_{\R^d} \hat f(\w) \prod_{i\in \vec u} \left(\e^{\im \omega_i x_i}- \hat{\mu}_i(-\omega_i)\right) \prod_{i \in \vec u^c} \hat{\mu}_i(-\omega_i) \d \w\\
&\quad=\frac{1}{(2\pi)^d}\int_{\R^d} \hat f(\w) E(\x,\w,\mu,\vec u) \d \w.
\end{align*}
This shows~\eqref{eq:anova_hat}. 
%It follows directly that
%\begin{align*}
%\hat{f_{\vec u}} (\w)
%&= \begin{cases} \hat{f}(\w) &\text{ if } \supp \w = \vec u\\
 %0& \text{ else,}
%\end{cases}
%\end{align*}
%since $\e^{\im \omega_i x_i}- \hat{\mu}_i(-\omega_i) = 0$ if $\omega_i = 0$ for some $i \in \vec u$.
%This finishes the proof.
\end{proof}

\end{document}
